\section{High modes of the tiling}
\label{sec:high-modes-of-the-tiling}

In the high modes we assign distinct odd labels first, then construct even rows supported on their common neighbors. We must retain positive mass in every even row and keep each column sum small enough that normalization permits Hall's theorem.

For an even \(v\) of patch \(i\) call its bulk and crossing neighbor sets \(L_v\) and \(H'_v\) respectively. Write \(i(b)\) for an odd neighbor's own patch, \(D_j=D_{\pi_j}\), \(I_x(y)=\mathbf1_G(x,y)\) (using the common color). Ratio weights are set to zero off their specified support/gates, requiring positive denominators on support.

\begin{lemma}[High direct assignment]\label{lem:high-direct-assignment}
In high direct mode, for all sufficiently large members of the counterexample sequence, there is an odd injection whose hit-supported even rows have mass at least one half and column sums at most one half. Consequently these rows admit distinct representatives.
\end{lemma}

\Needspace{5\baselineskip}
\begin{proof}\leavevmode\par
\paragraph{Mass estimate.}
In high direct mode define the unnormalized even row
\[
 Z_v(x)=\sigma_v(x)\prod_{b\sim v} I_x(y_b)/D_{i(b)}(x).
\]
First sample the odd labels independently from their \(\pi_{i(b)}\). We claim that \(\sum_xZ_v(x)<1/2\) has probability at most \(n^{-R}\).

Expose crossings first. After successful earlier filters, their normalized first law is at most \(4^{\ell_i}\sigma_v\), hence still has small width. Deep discrepancy \eqref{eq:source-9}, applied after conditioning the next second law on a putative exceptional set, puts its hit mass at \(1/2\pm O(b_*)\) except with probability exponentially small in a power of \(n\). The cleaned crossing denominators have the same bounds. Each normalized ratio filter therefore changes total mass by \(1+O(b_*)\), and the product is \(1+o(1)\), since \(\ell_i b_*=o(1)\).

Let \(\tau\) be the normalized first law after these filters. Its remaining bulk draws all have law \(\pi_i\). Its support retains the own-patch degree and clique conditions from Lemma~\ref{lem:stable-patch-cleaning}. The degree error is \(O(b_*)\), because \(g\le n^{\iota/2}\), and the positive surplus gives
\[
 D_i(x)^{-|L_v|}\le2^{|L_v|}e^{-.4g}.
\]
Here \(|L_v|=n-o(n)\). The support and crossing losses from allocation, together with \(N\ge2^n\) and \(s_i=10^{-3}g\), yield
\[
 \Delta_i\max_x\tau(x)D_i(x)^{-|L_v|}\le e^{-200s_i}.
\]
Apply Lemmas~\ref{lem:moderate-interaction-moment} and~\ref{lem:homogeneous-peeling} to this homogeneous bulk mass. Their even centered-moment bound also controls any fixed lower-tail threshold near \(1/2\). Since \(s_i>10^{-3}K_B\log n\), the peeling term and the crossing failures give the asserted \(n^{-R}\) bound.

\paragraph{Injection and column loads.}
Apply the clock lemma with \(B=5\), using these mass failures as predicates. Disjoint patches and the allocation counts give odd column sums \(o(1)\); atoms are exponentially small, each predicate reads \(n\) rows, and each row belongs to at most \(n\) predicates. The sampler therefore enforces every even-row mass bound.

Column loads require a separate estimate under this sampler. For fixed \(x\in X_i\), expand the \(n\)-th moment of the average of \(M_iZ_v(x)\) over even roles in patch \(i\). Its cap is \(2^ne^{-200s_i}\). Stars can intersect only at positions within distance 2, a fraction \(O(n^2)2^{-n+\ell_i}\); their cap contribution tends to zero because the high gain exceeds the polynomial and prefix losses. For the remaining positions the stars are disjoint. The clock comparison on their at most \(n^2\) queried rows restores independent labels, under which each row has mean \(M_i\sigma_v(x)=O(1)\). Lemma~\ref{lem:scattered-moments} gives a moment bound \(K^n\).

If this normalized average is \(A_x\), the actual column sum is
\(O(2^nM_i/N)A_x/M_i=O(A_x/C_n)\).
Markov and the union over at most \(N\le n2^n\) first labels bound the probability of any column sum above \(1/2\) by \(N(K'/C_n)^n=o(1)\). Thus one injection has both the enforced mass bounds and these column bounds. Normalize the hit-supported rows, increasing columns by at most a factor two, and apply Hall.
\end{proof}

\subsection{Cluster history alarms}
\label{sec:cluster-history-alarms}

Now use either high cluster mode, with the internal experiments and high profiles above. The stages are: raw primitive histories \(W\); avoidance of the local history failures defined below; a global test of the odd column loads; bin sampling conditional on each passing history; and label injection conditional on each successful bin history. These last samplers are normalized separately at their entering histories and do not reweight earlier stages. Let \(D_b^W(x)=\int I_x\,dp_{g(b)}^W\), \(g(b)\) the group of \(b\). Define \(J_v^0\subseteq X_i\) by
\[
 |D_b^W(x)-.5|\le2b_*\ (b\in L_v),\qquad
 1/2\le \prod_{b\in L_v}D_b^W(x)/D_i(x)\le2
\]
(requiring defined ratios); define \(J_v\) by further requiring \(|D_b^W(x)-.5|\le2b_*\) for crossings. Use
\begin{equation}
 Z_v(x)=\sigma_v(x)\mathbf1_{J_v}(x)
       \prod_{b\in L_v\cup H_v'} I_x(y_b)/D_b^W(x).             \label{eq:source-18}
\end{equation}
The row \(\sigma_v\) depends on \(W\) and the observed internal labels, while \(J_v^0\) depends only on the bulk histories. Those bulk inputs are in distinct own-patch slices different from \(v\)'s internal slice. Crossings are in other, mutually distinct patches. Neither of the next two history tests reads a crossing history; the further crossing removal will have a deterministic uniform-mass bound.

\paragraph{The product-of-degrees test.}
Along with \eqref{eq:source-15}, require the uniform \(X_i\)-mass of \((J_v^0)^c\) to be at most \(e^{-n^{.2}}\). Its raw failure probability is at most \(e^{-n^{.2}}\), as follows. Each law \(p_{g(b)}^W\) has small width for every history, so \eqref{eq:source-9} bounds its bulk degree-outlier fraction in \(X_i\) by \(e^{-\Omega(n)}\). Averaging in \(W\) and applying Markov shows that, outside an \(e^{-\Omega(n)}\) fraction of \(x\)'s, the probability of any such outlier is \(e^{-\Omega(n)}\).

For a remaining \(x\), the variables \(D_b^W(x)\), \(b\in L_v\), are independent by slices, with means \(D_i(x)\) by the high-profile identity. Clipping them to \(.5\pm2b_*\) changes those means negligibly. Their ranges are \(O(b_*)\), so concentration bounds a constant deviation of their sum by \(e^{-n^{.5}}\). Expanding the logarithms costs only \(O(nb_*^2)=O(n^{-.92})\). Thus their normalized product lies in \([1/2,2]\) except with this probability and the clipping exceptions. Integration over \(x\), then Markov at \(e^{-n^{.2}}\), proves the test bound. For each fixed crossing history, deep discrepancy also makes its extra removed \(X_i\)-mass \(e^{-\Omega(n)}\).

\paragraph{The large-interaction test.}
For \(\boldsymbol x=(x_1,\ldots,x_u)\), let \(M^W_{b,J}(\boldsymbol x)\) be the interactions of \eqref{eq:source-10} for the conditional bulk law \(p_{g(b)}^W\), with its own degree denominator. They are used only where all coordinates satisfy its degree gates. Define the positive function
\[
 \Phi_v^W(\boldsymbol x)=
 \sum_{I\subseteq[u]}\prod_{b\in L_v}
      \int\prod_{j\in I}\frac{I_{x_j}(y)}{D_i(x_j)}\,dp_{g(b)}^W(y)
\]
and the following nonnegative function of \(W\):
\[
 T_v(W)=\mathbb E_{\rm int\mid W}\!\left[
  \int\mathbf1_{(J_v^0)^u}(\boldsymbol x)
  \mathbf1_{\{\max_{b,J}|M^W_{b,J}(\boldsymbol x)|>2\xi\}}
       \Phi_v^W(\boldsymbol x)\,d\sigma_v^{\otimes u}(\boldsymbol x)
  \right].
\]
The integral is zero when \(\sigma_v=0\). Our second test is \(T_v(W)\le e^{-100s_i}\). It bounds an average over the internal reference experiment, not each realized internal posterior.

We prove the raw mean bound \(\mathbb ET_v\le e^{-300s_i}\) by comparison with the nominal interactions \(M_J^{\pi_i}\), which use the fixed law \(\pi_i\) and denominator \(D_i\).
If some nominal interaction exceeds \(\xi\), discard the conditional anomaly and degree indicators in this nonnegative integral. Averaging the independent bulk histories then replaces each \(p_{g(b)}^W\) by \(\pi_i\). Conditional on the internal data, the nonzero \(\sigma_v\) is supported on one cleaned corner support. Its widths and degrees fit the homogeneous estimates, and
\[
 \Delta_i\max_x\sigma_v(x)D_i(x)^{-|L_v|}\le e^{-450s_i}
\]
by \eqref{eq:source-12}, \eqref{eq:source-14}, the allocation slacks, and \(h_i+|L_v|\le n\). Lemma~\ref{lem:homogeneous-peeling} bounds this part.

On tuples with all nominal interactions at most \(\xi\), choose a bulk column witnessing a conditional interaction above \(2\xi\), and keep that column's degree gates. Freeze its history and all but one involved first coordinate. The raw internal experiment, and hence \(\sigma_v\) conditional on it, is independent of this bulk history. The gated one-free bound in Lemma~\ref{lem:deep-interaction-tails} therefore charges at most \(e^{-\alpha n/4}\) for the remaining coordinate. Meanwhile the other bulk histories integrate to their nominal laws; expanding each nominal factor bounds their product by \(e^{2^u\xi n}\). The singled-out column's factor is bounded by a constant depending on \(u\). Summing over columns and interactions, the choice of \(\xi\) makes this contribution exponentially small in \(n\), hence smaller than the required gain-scale bound. Together the two parts give the asserted raw mean.

The two test failures and the failure of \eqref{eq:source-15}, at each \(v\), thus have probabilities \(\le e^{-100s_i}\) by Markov and the high margins. Their \(W\)-scopes use only patch \(i\), at bounded complementary-word distance (outside internal axes) and inner distance \(<20\rho h_i\). Dependence degrees are at most \(n^{O(1)}\) times the volume of an inner ball of radius \(50\rho h_i\). Charges \(e^{-s_i}\) work for product-local-lemma conditioning, by the volume bound and \(s_i/\log n\to\infty\) uniformly here. Use this conditioned \(W\) law; dropping constraints touching polynomially many local consultations of the same sizes (possibly at different patches) costs \(1+o(1)\).

\begin{lemma}[Conditional high-cluster mass]\label{lem:high-cluster-mass}
Fix a history satisfying the two alarms above and \eqref{eq:source-15}. Given such \(W\), in product-by-groups reference odd sampling,
\begin{equation}
 \Pr(\textstyle\sum_x Z_v(x)<1/2\mid W)\le n^{-R}.                    \label{eq:source-19}
\end{equation}
\end{lemma}

\begin{proof}
By \eqref{eq:source-15}, the internal row \(\sigma_v\) is a probability except with probability \(\epsilon_i\). When it is nonzero, its mass removed by \(J_v\) is at most
\[
 \frac{M_i}{N}2^{h_i}e^{-500s_i}
       \bigl(e^{-n^{.2}}+e^{-\Omega(n)}\bigr)=o(1),
\]
using the first test, the deterministic crossing-removal bound, and \(h_i<n^\iota\). Now expose crossings. At each step the current normalized first law has small width, the new conditional law \(p_{g(b)}^W\) has small width, and its denominator is \(.5\pm2b_*\) on \(J_v\). The deep-discrepancy calculation from the direct case therefore makes each ratio filter change mass by \(1+O(b_*)\), outside exponentially small exceptions. At the end the retained mass is \(1+o(1)\), and its normalized law satisfies \(\tau\le5^{\ell_i+1}\sigma_v\) on \(J_v\).

The remaining bulk inputs are independent, each in a distinct group and slice. Apply \eqref{eq:source-10} with their laws \(p_{g(b)}^W\) and denominators \(D_b^W\). Lemma~\ref{lem:moderate-interaction-moment} bounds the small and moderate part of the centered moment by \(n^{-3R}\). For the large-interaction part, the domination of \(\tau\) and the bulk product-ratio gates give, after averaging the internal and crossing data with their preceding success indicators,
\[
 \text{large-interaction contribution}
       \le 5^{u(\ell_i+1)}2^u T_v(W).
\]
The factor \(2^u\) changes the conditional denominators to the nominal \(D_i\)'s in \(\Phi_v^W\). Since \(\ell_i\le s_i/(1000u)\), the alarm bounds this expression by \(n^{-3R}\). Even centered-power Markov, adding the preceding exceptions, gives \eqref{eq:source-19}.
\end{proof}

\subsection{Bin and injection stages}
\label{sec:bin-and-injection-stages}

We first prove that, with high probability under the history-avoidance law, every odd-label column sum is less than a fixed \(\theta_*>0\), chosen sufficiently small compared with the clock constant \(\theta_0\). For a label in patch \(i\), take the \(n\)-th moment of the average of \(M_ip_{g(b)}^W(y)\) over its odd roles. Equation~\eqref{eq:source-13} bounds each term by \(8e^{2k_iT_i}\). Treat two roles in the same slice as near; their fraction is at most \(2^{-n+o(n)}\), so the repeated-slice contribution in Lemma~\ref{lem:scattered-moments} tends to zero. For retained roles in distinct slices, dropping avoidance constraints touching their local consultations costs \(1+o(1)\) and restores independent raw histories. Their means are \(M_i\pi_i(y)\le11\) by \eqref{eq:source-16}. The resulting moment is \(K^n\), and the patch-count conversion used in the direct case makes every actual column sum \(o(1)\) with high probability. Denote this global load success by \(\mathcal H\). It is not an additional local conditioning of \(W\). The following samplers are used only on \(\mathcal H\); off it extend the experiment arbitrarily and set all tested weights to zero. Thus later moment estimates retain \(\mathbf1_{\mathcal H}\) until the conditional sampling comparisons have been applied.

\begin{proposition}[High-cluster sampling stages]\label{prop:high-cluster-sampling}
Fix an admissible \(W\)-history with the preceding column-load success. There is a conditional bin sampler satisfying the star mass tests, the small-bin distinctness tests when relevant, and the prescribed label-column capacity. On the stated local scopes it has upper comparison with independent group choices. For every successful bin history there is then an odd injection satisfying all even-row mass gates, with the label-stage comparisons described below. These stages are conditional samplers; completing a stage does not reweight its incoming history.
\end{proposition}

\Needspace{5\baselineskip}
\begin{proof}\leavevmode\par
\paragraph{Bin predicates.}
Given such prehistory, start the bin stage with independent group choices by \(q_g^W\) (bin identities below refer to physical bins, with their particular further masked/filtered laws as in each group). Require per \(v\) that conditional on bins its independent-role label probability of failure in \eqref{eq:source-19} is \(\le n^{-R/2}\); exception costs \(\le n^{-R/2}\) by Markov. In small-bin mode also require no bin chosen twice by distinct groups in its star, with exponentially small exception by \eqref{eq:source-13}. These are star constraints on bin choices. Finally require given-bin label column sums \(\le\theta_0\).

\paragraph{Capacity certificates.}
Fix a physical bin \(D\) of size \(d=d_i\), a label \(y\in D\), and the entering history \(W\). If group \(g\) chooses this bin, its \(h_i\) roles contribute
\[
 a_g(y):=h_iU_g^W(D)(y)\le d^{-1/2}
\]
to the column, by \eqref{eq:source-13} and the scale choices. Put groups with \(a_g(y)\in(\lambda_j/2,\lambda_j]\) in bucket \(j\), where \(\lambda_j=d^{-1/2}2^{-j}\). Its expected load under independent bin choices is
\(m_j=\sum_{g\text{ in bucket }j}q_g^W(D)a_g(y)\).
For
\[
 l_j=\left\lfloor(K'm_j+c'2^{-j/2})/\lambda_j\right\rfloor+1,
\]
forbid every event that a specified \(l_j\)-subset of this bucket all choose \(D\). Avoidance leaves at most \(l_j-1\) chosen groups in the bucket, so its load is at most \(K'm_j+c'2^{-j/2}\). First choose \(K'\) large, then \(c',\theta_*\) sufficiently small. Since \(\sum_jm_j<\theta_*\), the sum of these allowed loads is at most \(\theta_0\).

Give such a certificate charge \(2^{l_j}\) times its product probability. To sum charges touching a fixed group, first pin its bin choice to \(D\). For each \(y\in D\), that group belongs to only one bucket. In that bucket put \(z=\sum_gq_g^W(D)\le2m_j/\lambda_j\). The chosen constants and the additive allowance ensure
\[
 l_j-1\ge20z,\qquad l_j-1\gtrsim d^{1/2}2^{j/2}.
\]
Summing over subsets containing the pinned group costs at most
\[
 2^{l_j}\frac{z^{l_j-1}}{(l_j-1)!}
       \le e^{-\Omega(d^{.4})}.
\]
The union over at most \(d\) columns still has this negligible bound; averaging the pinned target gives the unconditional touching bound. Here \(q>(\log n)^{c_q}\), so these errors are smaller than every inverse polynomial in \(n\).

Use charge \(n^{-5P}\) for each star constraint. At most \(O(nh_i)\) of them touch one group, and their raw probabilities have much stronger decay. Thus the total charge touching any group is \(\delta_n=o(n^{-4P})\). For a certificate involving \(l\) groups, the reciprocal neighbor product is at most \(\exp(O(l\delta_n))\le2^l\), which is paid by its chosen charge inflation. The polynomial scopes of star constraints likewise have reciprocal neighbor product \(1+o(1)\). Lemma~\ref{lem:conditional-avoidance} now supplies all bin requirements. Dropping constraints touching \(O(n^3)\) queried groups costs \(1+o(1)\), giving the asserted comparison with independent bin choices.

\paragraph{Conditional label assignment.}
In high-large-bin mode individual laws conditional on successful bins now meet the clock lemma (atoms \(\le2\exp(1.5k_iT_i)/d_i\) uniformly superpolynomially small). Draw an odd injection avoiding every mass failure. In high-small-bin mode instead start independently per assigned bin with \eqref{eq:source-17} for its individual laws (maximum atom \(\le d_i^{-.95}\) by \eqref{eq:source-13}). Here \(d_i\le n^{o(1)}\) and \(h_i=o(d_i^{.025})\) uniformly. In each star a queried bin belongs to just one participating group by bin-stage success; multi-observations in it are possible only for internal incidences. Using exact singleton comparisons and \eqref{eq:source-17} their overall joint upper-error factor is \(\le\exp(d_i^{-.04}h_i)\), so failure still costs \(\le2n^{-R/2}\). Each bin serves at most \(d_i\) roles by the column-sum requirement, hence bin-variable dependence degrees are \(n^{2+o(1)}\) or less. Condition to avoid all these failures by the local lemma (charges e.g. \(n^{-5P}\)), retaining upper comparison to the preconditioning bin-wise sampling at \(1+o(1)\) cost on \(O(n^2)\) queried bins. Thus either procedure gives an odd injection with the mass gates, without further reweighting earlier histories.
\end{proof}

\subsection{High cluster load transfer}
\label{sec:high-cluster-load-transfer}

\begin{proposition}[High-cluster column loads]\label{prop:high-cluster-loads}
Let \(\mathcal E_i\) be the even positions of patch \(i\), and let \(\mathcal H\) denote the prerequisite \(W\)-load success. Under the conditioned history law and the conditional samplers of Proposition~\ref{prop:high-cluster-sampling}, uniformly for \(x\) in the cleaned envelope,
\[
 \mathbb E\!\left[\mathbf1_{\mathcal H}
 \left(\frac{1}{|\mathcal E_i|}\sum_{v\in\mathcal E_i}M_iZ_v(x)\right)^n\right]
 \le K^n .
\]
Together with the mass gates, the patch counts and divergence of the ambient ratio, this gives an embedding in each high cluster mode.
\end{proposition}

\Needspace{5\baselineskip}
\begin{proof}\leavevmode\par
\paragraph{Caps and geometric separation.}
Fix \(x\) in patch \(i\)'s cleaned envelope; only its even patch positions can use it. We bound \(n\)-th moments of the patch average of \(M_i Z_v(x)\), with indicator of the prerequisite \(W\)-load success. On successes use both
\begin{equation}
 M_i Z_v(x)\le 2^n e^{-200s_i},\qquad
 Z_v(x)\le 4\,\sigma_v(x)\prod_{b\in L_v\cup H_v'} I_x(y_b)/D_{i(b)}(x).     \label{eq:source-20}
\end{equation}
Indeed off \(J_v\) the weight is zero; on it the bulk ratio gate and crossing degree bounds for both kernels (\(\ell_i b_*=o(1)\)) give the product comparison. Its individual factors are \(\le3\) on this envelope. Also use \eqref{eq:source-12}, \eqref{eq:source-14} with \(10q^{C_b}\le s_i/(100u)\) for the cap.

There are three possible sources of repeated information in this moment: intersecting geometric core scopes, crossing scopes shared by different stars, and, for small bins, distinct groups choosing the same physical bin. We handle them in that order.

Expand by independent uniform even patch indices (positions, ordered by draw number). Call a position's core scopes those for \(\sigma_v\) and its internal and bulk samples/kernels: the needed own-slice and nearby own-patch \(W\) consultations, groups and odd roles. Its core bin-choice variables are just the groups actually sampled for its internal and bulk labels; even the counterfactual calculations for \(\sigma_v\) use the observed internal labels and \(W\), not bin draws elsewhere. Remove by the cap each index core-near an earlier index (earlier removals included), using proximity of complementary-word distance \(\le4\) \textbf{and} inner distance \(\le100\rho h_i\). This covers core intersections, and partner fraction per fixed index is at most \(2^{-n}e^{.01s_i}\) by the inner volume, prefix and high slacks. On retained indices form a graph by full distance \(\le n^{2\iota}\), partner fraction \(\le2^{-n+o(n)}\). It covers any crossing-scope intersections between them (all internal dimensions \(<n^\iota\)); crossing consultations themselves are never in patch \(i\), and within a star are in different patches. When using the second bound in \eqref{eq:source-20} below, drop all crossing factors at nonisolated indices of this graph. If this graph has rank \(j\), meaning its number of vertices minus its number of connected components, at most \(2j\) vertices are nonisolated. Deleting their crossing factors therefore costs at most \(3^{2j\ell_i}\).

\paragraph{Label and bin transfers.}
For large bins transfer the retained indices' needed data in that product bound to independent labels given bins by the clock bound (at most \(n^2\) distinct queries), then use independent-bin upper comparison on the involved groups.

For small bins we need the following extra splice before these transfers. Given the bin choices, remove by the cap each geometrically retained index whose core groups' bin choices intersect the core choices of an earlier geometrically retained index (even one that these bin repeats remove). In using \eqref{eq:source-20} thus take the cap on all indicated removals, using the product bound and crossing-graph drops only on the remaining indices. For a core-bin removal keep its necessary repeat condition; against fixed earlier choices under independent bin sampling given \(W\) this costs per index
\[
 n^{O(1)}\,4e^{2k_iT_i}d_i/M_i \le 2^{-n}\exp(o(s_i)),
\]
by distinctness of group variables across the geometrically retained cores and the high scale bounds. Among the remaining indices' crossing queries not excised by the graph, use a fixed order and drop by factor 3 each hit factor whose bin repeats an earlier bin of these queries. Their group variables are all distinct and outside patch \(i\); the necessary repeat condition costs \(\le2^{-n+o(n)}\) per drop against earlier choices under independent bins.

\paragraph{Small-bin removal patterns.}
Fix the position tuple and then one outcome of the ordered removal procedure just described: core-near rows are capped first; among the remaining rows, later core-bin repeats are capped; among crossing queries still allowed by the crossing graph, later bin repeats lose only their hit factor. Fix the resulting removal masks before comparing probability laws.

Let \(\mathcal R\) be the rows whose factors remain, and \(\mathcal B_v\) their remaining bulk and crossing queries. After paying the row caps and the factors 3, the nonnegative integrand has the form
\[
 \mathbf1_{\mathrm{rep}}
 \prod_{v\in\mathcal R}
   \left[M_i\sigma_v(x)
      \prod_{b\in\mathcal B_v}\frac{I_x(y_b)}{D_{i(b)}(x)}\right].
\]
Here \(\mathbf1_{\mathrm{rep}}\) retains the necessary earlier-bin repeat for every removal. Initially one also retains the mask-consistency and successful-history indicators. The rule for \(\sigma_v\) uses its observed internal labels and local \(W\); its candidate calculations integrate hypothetical bins and never consult other actual bin choices. Consequently the displayed product reads actual bins only through its remaining label queries.

First, conditional on a successful bin history, remove touching label-stage constraints and apply \eqref{eq:source-17}. Different remaining groups use distinct bins; several role queries can share a bin only within one internal group. Exact singleton marginals handle the other queries. The total comparison cost is
\[
 (1+o(1))\exp(d_i^{-.04}h_i n)\le O(1)^n.
\]
Integrating the resulting independent role labels leaves a nonnegative function only of the kept bin choices and local \(W\), multiplied by the repeat indicators.

Next apply the bin-stage upper comparison to every group read by this function or a repeat condition. Its independent-role formula extends nonnegatively off the bin gates. Thus, after this comparison, discard the other bin gates and mask-consistency restrictions, retaining just the necessary earlier repeats.

At this point the bin variables have their independent raw laws given \(W\). Integrate removed variables in reverse order. Earlier removed choices can be held fixed while charging the necessary repeat of a later one; the surviving label integral does not read that later variable. A core-row repeat costs at most \(2^{-n}e^{o(s_i)}\), paying its cap \(2^ne^{-200s_i}\). A crossing repeat costs at most \(2^{-n+o(n)}\), paying its factor 3. The two types concern disjoint patches. There are at most \(2^n\) core masks, while the crossing-repeat costs sum to \(1+o(1)\) over masks on at most \(n^2\) queries. All these bounds are uniform in the kept choices and \(W\).

\paragraph{Restoring the reference product.}
Reverse integration has therefore used independence only after both the label and bin comparisons, not under the avoidance-conditioned bin law. The prerequisite \(W\)-load indicator has remained in place to license those comparisons.

After these transfers, integrate kept independent bin choices and reference label draws. Keep only the internal reference weights and used external factors in the product test; in particular nonlocal \(W\)-load restrictions can now be dropped. Use the \(W\) upper comparison on the kept local scopes. In unrestricted sampling they give independent core experiments across positions, internal data independent of the used external inputs, and independence across those individual bulk and remaining crossing inputs. This follows by core separation and removal of crossing overlaps as above and the slice/patch structure within each star. Thus each \(I_x(y_b)/D_{i(b)}(x)\) left integrates to 1 by profile definition, while \(M_i\sigma_v(x)\) on internal reference data has mean \(\le K\) by \eqref{eq:source-14}. The product-mean costs here are \(O(1)^n\).

\paragraph{Moment summation and Hall.}
Finally enumerate geometric removal masks (at most \(2^n\)); in independent uniform-position sampling conditional on the values at a specified mask's retained indices, enforcing just necessary earlier proximities for its removed indices and summing those variables backwards costs per removal at most \(n\,2^{-n}e^{.01s_i}\). The graph cost on retained positions (before bin splices) for rank \(j\) needs a forest of \(j\) edges paying fraction \((2^{-n+o(n)})^j\) under unrestricted uniform-position sampling, with at most \(n^{2j}\) choices. This absorbs \(3^{2j\ell_i}\). Thus all normalized moments in question are \(O(1)^n\) uniformly. For the normalized column average \(A_x=|\mathcal E_i|^{-1}\sum_v M_iZ_v(x)\), the actual column sum is
\[
 \sum_{v\in\mathcal E_i}Z_v(x)
   =\frac{|\mathcal E_i|}{M_i}A_x=O(C_n^{-1})A_x.
\]
The moment bound with \(\mathbf1_{\mathcal H}\), Markov, and the union over at most \(N\le n2^n\) labels give failure probability at most \(N(K'/C_n)^n=o(1)\). Meanwhile \(\Pr(\mathcal H)=1-o(1)\), and both conditional samplers enforce their mass gates at every passing entering history. Their intersection with column sums at most \(1/2\) is therefore nonempty. Normalizing row masses, which are at least \(1/2\), and applying Hall gives the cube. This excludes the high modes of the tiling.
\end{proof}
